\documentclass[11pt]{article}
% Standalone preamble for web-only supplementary notes.
\usepackage[T1]{fontenc}
\usepackage{amsmath,amssymb,amsthm}
\usepackage{hyperref}
\usepackage{cleveref}
\DeclareMathOperator{\Map}{Map}
\DeclareMathOperator{\Fun}{Fun}
\newcommand{\iso}{\xrightarrow{\sim}}
\newtheorem{lemma}{Supplementary Lemma}
\newtheorem{corollary}[lemma]{Supplementary Corollary}
\crefname{lemma}{Supplementary Lemma}{Supplementary Lemmas}
\Crefname{lemma}{Supplementary Lemma}{Supplementary Lemmas}
\crefname{corollary}{Supplementary Corollary}{Supplementary Corollaries}
\Crefname{corollary}{Supplementary Corollary}{Supplementary Corollaries}

\begin{document}
\section*{Lifting prescribed identifications through localization}

When using the universal property of a
\href{localizations.html}{localization}, it can matter which identification
is obtained after restriction. Merely identifying the two extended functors
does not compute the image of that identification. This note records the
stronger calculation used in the formalization.
Assume the Chapter 1 mapping-anima and pullback structures.

\begin{lemma}
\label[lemma]{lem:Supplementary_Restriction_Image}
Let $r\colon B\to C$ and suppose
$r^*\colon\Map(C,D)\to\Map(B,D)$ is an embedding.
For functors $f,g\colon C\to D$ and a specified identification
$\alpha\colon fr=gr$, there is an identification $\widetilde\alpha\colon f=g$
with a specified higher identification
\[
 \widetilde\alpha|_r=\alpha.
\]
\end{lemma}
\begin{proof}[Proof idea]
Name $f$ and $g$ in $\Map(C,D)$. Move $\alpha$ across the comparisons
between restriction of their names and names of their restrictions.
The embedding $r^*$ induces an equivalence of identification animae,
so the resulting identification lifts with its prescribed image.
Decode the lift and transport its endpoints back to $f,g$.

The computation law for this decoding, together with its compatibility
with restriction, proves the displayed image formula. The final endpoint
transports are essential: the lifted identification initially has the
decoded named functors as its endpoints. See also
\href{supplementary-decoding-and-restriction.html}{Decoding mapping animae
and restriction}.
\end{proof}

\begin{corollary}
\label[corollary]{cor:Supplementary_Localization_Image}
Let $l\colon C\to L$ be a localization at a collection $W$, in the
setting of the book's localization definition. For every target $D$,
restriction $\Map(L,D)\to\Map(C,D)$ is an embedding.
Consequently every specified identification between the restrictions of
two functors $L\to D$ lifts as in
\Cref{lem:Supplementary_Restriction_Image}.
\end{corollary}
\begin{proof}[Proof idea]
Restriction factors through the equivalence
$\Map(L,D)\iso\Map^W(C,D)$ and the inclusion of this subanima in
$\Map(C,D)$. A composite of an equivalence and an embedding is an
embedding. The ordinary restriction map need not be an equivalence:
its image consists of functors inverting $W$.
\end{proof}

\begin{lemma}
\label[lemma]{lem:Supplementary_Restriction_Higher_Reflection}
If $r^*\colon\Map(C,D)\to\Map(B,D)$ is an equivalence, then for
$\alpha,\beta\colon f=g$, an identification
$\alpha|_r=\beta|_r$ induces an identification $\alpha=\beta$.
\end{lemma}
\begin{proof}[Proof idea]
Encode $\alpha$ and $\beta$ as identifications between the names of
$f$ and $g$. Naturality of decoding turns the hypothesis into an
identification of their restricted encodings. Reflect through the
induced equivalence of identification animae, decode, and cancel the
endpoint comparisons. The equivalence hypothesis is the one used by
this higher-reflection declaration.
\end{proof}

These two levels of calculation occur in
\href{geometric-realizations.html}{preservation of pushouts by geometric
realization}. When reconstructing a cocone through restriction equivalences,
first lift the legs and then lift the specified matching with
\Cref{lem:Supplementary_Restriction_Image}. For a comparison of cocones,
lift both leg comparisons with prescribed images. Its matching equation
can then be checked after restriction and reflected by
\Cref{lem:Supplementary_Restriction_Higher_Reflection}. Thus the comparison
retains the whole cocone. The same distinction is useful when reading the
\href{chapter-3-exercises.html}{uniqueness exercise for localizations}.
\end{document}
